\section{D. Quasi-unipotence, integral one-parameter groups, and the presentation request}
\label{sec:D}

We distinguish a homomorphism on a cyclic group from its germ on finite-index
subgroups.  This distinction is essential in \S52.3, fols.~752'--753,
C4 scans 210--211.  We also correct a location error in the development
prompt: C4 scan 231 is already \S53.3, fol.~763.  The presentations begun
there are proved on scans 234--237 and concern a central covering of a
modular group; they are not presentations of a nilpotent Lie quotient.
We prove them in Part E and explain below why their nilpotent tests have
limited information.

The predecessor arguments \citep[Part E, E.3, E.6, E.10--E.11]{LM1Return2026}
have been reopened and checked; the needed proofs are reproduced below.
The nonlinear argument is completed in D.17--D.21, and the actual
smooth-model question is resolved with its exact scope in D.22--D.23.

\subsection{What a cyclic germ determines}

\begin{statement}{D.1}{Proved}{The rational one-parameter subgroup of a quasi-unipotent germ}
Let $G\hookrightarrow\mathrm{GL}_{r,\mathbb Q}$ be a closed algebraic
subgroup, and let $g\in G(\mathbb Q)$ be quasi-unipotent.  Choose $n>0$
such that $g^n$ is unipotent.  There is a unique homomorphism of
$\mathbb Q$-group schemes
\[
 u_g:\mathbb G_{a,\mathbb Q}\longrightarrow G,
 \qquad u_g(t)=\exp\bigl(t\log(g^n)/n\bigr)
\]
whose restriction to $n\mathbb Z$ agrees with $k\mapsto g^k$.
It is independent of the choice of $n$.  It is a monomorphism if it is
nontrivial.  In general $u_g(1)\ne g$.
\end{statement}
\begin{proof}
Set $N=g^n-I$.  Both logarithm and exponential are finite polynomials
in the nilpotent matrix $N$.  The identity
$\exp((s+t)X)=\exp(sX)\exp(tX)$ for $X=\log(g^n)/n$ proves that
the displayed matrix polynomial is a homomorphism to $\mathrm{GL}_r$.
For every $k\in\mathbb Z$ its value at $nk$ is $g^{nk}\in G$.
Every equation defining $G$ therefore vanishes after substitution:
a rational polynomial vanishing at the infinite set $n\mathbb Z$
is zero.  This proves factorization through $G$.

For uniqueness, if $P(t)$ is any polynomial matrix homomorphism,
differentiating $P(t+s)=P(t)P(s)$ at $s=0$ gives
$P'(t)=P(t)P'(0)$.  Coefficient comparison yields
$P(t)=\exp(tP'(0))$.  Its value at $n$ fixes its logarithm, hence
$P'(0)=X$.  Taking a common multiple of two choices proves independence.
If $X\ne0$, the kernel is a proper algebraic subgroup of
$\mathbb G_a$ in characteristic zero.  Such a subgroup is trivial:
its Lie algebra is either $0$ or the full one-dimensional Lie algebra,
and a finite subgroup has no nonzero element in characteristic zero.
The differential is injective as well, so the kernel group scheme is
trivial.  Finally $G=\mathbb G_m$, $g=-1$, gives $u_g=1$ and
$u_g(1)=1\ne-1$.  Thus scan 211 is correct as a statement about
the germ, exactly as its phrase ``pas seulement un germe'' requires;
replacing germ agreement by agreement on the whole original cyclic
group is false.
\end{proof}

\begin{statement}{D.2}{Proved}{Integral divided powers and positive powers}
Let $A$ be a torsion-free $\mathbb Z$-algebra and let
$G\hookrightarrow\mathrm{GL}_{r,A}$ be a closed subgroup scheme.
For $g\in G(A)$ with $g-I$ nilpotent, put $X=\log g$ over
$A\otimes\mathbb Q$.  A homomorphism
$\mathbb G_{a,A}\to G$ taking $1$ to $g$ exists exactly when
\[
 X^j/j!\in M_r(A)\qquad(j\geq0);
\]
it is then unique.  More generally, if $A$ is any commutative ring,
$G$ is a closed finitely presented subgroup of $\mathrm{GL}_{r,A}$,
and $g$ is fibrewise quasi-unipotent in Grothendieck's sense, then
some positive power $g^m$ is the value at $1$ of a homomorphism
$\mathbb G_{a,A}\to G$.
\end{statement}
\begin{proof}
The rational proof in D.1, with $n=1$, applies over any
$\mathbb Q$-algebra: nilpotence makes the logarithm finite, and
Vandermonde determinants on finitely many integer evaluations are
units.  It forces the displayed coefficients.  Since
$A\hookrightarrow A\otimes\mathbb Q$ is injective, their integrality
is necessary, and if they are integral the addition identity and
the defining equations of $G$ descend by this injectivity.  This
proves the first assertion.

For the second, in every residue field a positive power of $g$ has
nilpotent difference from $I$.  The constructibly clopen subsets
defined by the entries of $(g^a-I)^r$, as $a$ varies over positive
integers, cover $\operatorname{Spec}A$.  Compactness of the
constructible topology gives a finite subcover.  Take a common
multiple $a$ of its exponents.  The coefficients of the characteristic
polynomial of $N=g^a-I$, apart from the leading coefficient, are
nilpotent.  They generate a nilpotent ideal because there are finitely
many of them.  Cayley--Hamilton then gives $N^b=0$ for some $b$.
For $d=(b-1)!$ the polynomials $\binom{dt}{j}$, $j<b$, lie in
$\mathbb Z[t]$: each nonconstant coefficient of the numerator is
divisible by $d$, hence by $j!$.  Consequently
\[
 P(t)=\sum_{j=0}^{b-1}\binom{dt}{j}N^j
\]
defines an additive one-parameter group in $\mathrm{GL}_{r,A}$,
by the binomial identity in the universal ring
$\mathbb Z[N]/(N^b)$, and $P(k)=g^{adk}$ at each integer $k$.

It remains to enforce the equations of $G$; integer values alone
are insufficient over a ring with torsion.  If a polynomial
$F(t)\in A[t]$ vanishes at all integers, evaluation on
$0,\ldots,\deg F$ and multiplication by the adjugate of the
Vandermonde matrix show that some nonzero integer annihilates
every coefficient.  If also $F(0)=0$, replacing $t$ by $Dt$ for
that integer makes $F$ identically zero.  Apply this to a finite
set of defining equations of $G$, and use a common dilation $D$.
Then $P(Dt)$ factors through $G$ and takes $1$ to $g^{adD}$.
This proves, with its linearity and finite-presentation hypotheses,
the envelope argument used in \S52.3, fol.~754, scan 213.
\end{proof}

\begin{statement}{D.3}{Refuted}{Fibrewise unipotence does not supply an integral parameter}
The assertion that every fibrewise unipotent integral matrix is the
value at $1$ of an integral additive one-parameter subgroup is false.
Let $N\in M_5(\mathbb Z)$ have superdiagonal $(1,2,3,2)$ and zero
other entries.  Then $g=I+N$ is fibrewise unipotent in the
manuscript's one-parameter sense, but no homomorphism
$\mathbb G_{a,\mathbb Z}\to\mathrm{GL}_{5,\mathbb Z}$ takes $1$ to $g$.
The exact replacement is D.2, not fibrewise unipotence alone.
\end{statement}
\begin{proof}
We have $N^2=0$ modulo $2$, $N^3=0$ modulo $3$, and $N^5=0$.
In characteristic $p$, a nilpotent matrix $\overline N$ with
$\overline N^p=0$ has a truncated logarithm and exponential with
all denominators invertible, giving the required homomorphism.
This treats $p=2,3$ and every $p\geq5$; characteristic zero is D.1.
Direct multiplication gives
\[
 X=\log(I+N)=
 \begin{pmatrix}
 0&1&-1&2&-3\\0&0&2&-3&4\\0&0&0&3&-3\\
 0&0&0&0&2\\0&0&0&0&0
 \end{pmatrix},\qquad X^4/4!=\tfrac12 E_{15}.
\]
D.2 supplies the contradiction.  The exact arithmetic is independently
recomputed in \texttt{code/d\_checks.py}; its JSON and text outputs
record the matrices and every divided-power denominator.
\end{proof}

\begin{statement}{D.4}{Refuted}{Smooth closure and the primitive lattice are distinct conditions}
Smoothness of the closure of a rational additive subgroup does not
imply that the map on any positive-index sublattice is a monomorphism,
nor that the primitive integral parameter comes from the original
discrete group.  For
\[
 \Pi_0=\mathbb Z,\qquad \varphi(k)=2k\in\mathbb G_a(\mathbb Z),
\]
the closure is $\mathbb G_{a,\mathbb Z}$, but the map from $n\mathbb Z$
has coefficient $2n$ and is never a monomorphism over $\mathbb Z$.
The saturated positive parameter is $t\mapsto t$; it does not
factor through $\varphi(\Pi_0)$ on integral points.
\end{statement}
\begin{proof}
The rational image is all $\mathbb G_{a,\mathbb Q}$, whose schematic
closure is $\mathbb G_{a,\mathbb Z}$.  The lattice $n\mathbb Z$
has generator $n$, so its additive group scheme maps by $t\mapsto2nt$.
Its fibre at a prime dividing $2n$ is the zero homomorphism, hence
has nontrivial kernel.  The only automorphisms of
$\mathbb G_{a,\mathbb Z}$ are $t\mapsto\pm t$; requiring a positive
rational scaling selects $t\mapsto t$.  Its value at $1$ is odd
and therefore outside $2\mathbb Z$.  This isolates two assertions
in scans 214--215, fols.~754'--755: the suggested sufficiency for
an index $n$ is false as stated, whereas the extra factorization
hypothesis (55) is genuinely extra.  For a chosen integral copy of
$\mathbb G_a$, the strongest elementary lattice criterion is that
the coefficient of the parameter be a unit; at a prime $p$ it must
be a $p$-adic unit.
\end{proof}

\begin{statement}{D.5}{Proved}{The Lie-spanning condition has a precise smooth form}
Let $H/\mathbb Z$ be a smooth affine group scheme of finite presentation
with geometrically connected fibres.  Suppose finitely many
$u_i:\mathbb G_a\to H$ have differentials spanning
$\operatorname{Lie}(H)$ over $\mathbb Z$.  Then their images generate
$H$ as an etale sheaf of groups.  The assertion also holds after
adjoining finitely many conjugate additive subgroups.  For a smooth
group with finite constant component group, one must additionally
supply sections meeting every component.
\end{statement}
\begin{proof}
The multiplication map from a finite product of the additive groups
to $H$ has surjective differential at the identity section.  Since
source and target are smooth, it is smooth on a neighbourhood of
that section.  Its image contains an open neighbourhood $V$ of the
identity, and every point of $V$ lifts etale locally to the source:
a smooth surjection admits sections after etale localization.
The subgroup sheaf generated by $V$ is represented by an open subgroup;
one may construct it as the union of the images of finite products
of $V\cup V^{-1}$.  In every geometric fibre this open subgroup is
the whole connected group, since its cosets are open and its complement
is a union of open cosets.  Therefore the subgroup sheaf is all of $H$.
Apply the same argument to $H^0$ when components are present and
then use the specified representatives.  This proves the usable
version of the implication on \S52.3, fols.~755'--756,
scans 216--217; it does not justify dropping smoothness or the
condition on components.
\end{proof}

\subsection{Two specified Malcev quotients and their integral matrices}

For this computation, and only for it, the two levels are
$F_r/\gamma_2F_r$ and $F_r/\gamma_3F_r$ followed by rational Malcev
completion.  They are auxiliary choices, not a claim that \S52
selects a canonical first and second approximation.  For integral
questions we specify the faithful truncated-Magnus linear model
\[
 T_{\leq d,\mathbb Z}
   =\mathbb Z\langle X_1,\ldots,X_r\rangle/(\text{words of length}>d),
 \qquad x_i\longmapsto1+X_i.
\]
The induced substitution action of a free-group automorphism is an
integral matrix in the word basis.  Over $\mathbb Q$ its restriction
to the Malcev subgroup gives the action on the corresponding
nilpotent quotient.

\begin{statement}{D.6}{Proved}{Torus and four-puncture twists at the two levels}
On $F_2=\langle x,y\rangle$, set
\[
 A(x)=x,\quad A(y)=yx,\qquad
 B(x)=xy^{-1},\quad B(y)=y,
 \qquad c=[x,y].
\]
Both fix $c$ and satisfy $ABA=BAB$.  On the free class-two rational
Lie algebra with $X=\log x$, $Y=\log y$, $C=[X,Y]$, the logarithm
$D_A$ of $A$ is
\[
 D_A(X)=0,\qquad D_A(Y)=X-\tfrac12C,\qquad D_A(C)=0.
\]
Thus $A^t=\exp(tD_A)$ is the unique rational additive parameter.
In the specified Magnus lattice it is integral at level one;
at level two its divided-power denominators are $1,2,2$, so the
least positive integral power is $A^2$.  At $3$ no dilation is needed.

For a sphere with four labelled punctures use
$F_3=\langle x_1,x_2,x_3\rangle$, the fourth loop
$(x_1x_2x_3)^{-1}$, and $p=x_1x_2$.  A separating Dehn twist is
\[
 T(x_1)=px_1p^{-1},\quad T(x_2)=px_2p^{-1},\quad T(x_3)=x_3.
\]
It is trivial at level one.  At level two its logarithm sends
$X_1$ to $-[X_1,X_2]$, $X_2$ to $[X_1,X_2]$, and $X_3$ to $0$.
Its square is zero as a linear operator, so its additive parameter
is integral in the Magnus lattice at both $2$ and $3$.
\end{statement}
\begin{proof}
Reduction of words gives $[x,yx]=[xy^{-1},y]=c$ and $ABA=BAB$.
The class-two Baker--Campbell--Hausdorff formula gives
$\log(yx)=Y+X+[Y,X]/2$, proving the first Lie formula.  In the
degree-two Magnus algebra the image of $Y$ is $Y+X+YX$.
If $N=A-I$ on this seven-dimensional word module, then
$N^3=0$ and $N^2(Y)=X^2$.  Consequently
\[
 \log A(Y)=X+YX-\tfrac12X^2,
 \qquad (\log A)^2(Y)/2=\tfrac12X^2.
\]
All remaining entries have denominators dividing $2$, as substitution
on $1,X,Y,X^2,XY,YX,Y^2$ verifies.  D.2 rules out an integral
parameter for $A$ and supplies one for $A^2$.  Denominator $2$ is
invertible over $\mathbb Z_{(3)}$.

For $T$, conjugation by $p$ has degree-two increment
$[X_1+X_2,X_i]$ on $X_i$, $i=1,2$.  It fixes $X_3$ and the
degree-two layer; its increment squares to zero.  Thus
$T^t=I+t(T-I)$, integrally.  The word identity
$T(x_1x_2)=x_1x_2$ also proves exact preservation of the fourth
loop.  All matrices, rational logarithms, and denominator tests are
computed by \texttt{d\_checks.py}; no floating-point rank test is used.
In the second Magnus representation the reduction of $A$ modulo $2$
has order $4$, so it is quasi-unipotent, but not itself a value on
$\mathbb G_a$ in that fibre.  The rank-five fibrewise-unipotent
counterexample D.3 concerns the stronger, undilated hypothesis.
The integral conclusion is explicitly about this lattice.  Changing
to a different group-scheme model requires a lattice comparison.
\end{proof}

\begin{statement}{D.7}{Proved}{What the requested class-four presentation tests detect}
The assignments $a\mapsto A$, $b\mapsto B$ define an action of
$B_3=\langle a,b\mid aba=bab\rangle$ on $F_2$.  The program verifies
the braid relation, the reflection relations of E.3, and
\[
 (AB)^6=\operatorname{Int}(c^{-1})
\]
on free nilpotent quotients of every class $1\leq d\leq4$.
These tests follow from identities of reduced free words, also
checked by the program.  Conversely, every nilpotent quotient of
$B_3$ is cyclic; testing the modular-cover presentation solely
in its own nilpotent quotients cannot establish it.
\end{statement}
\begin{proof}
The first statements are exact substitution identities in a free
group, so they descend along each characteristic quotient map
$F_2\to F_2/\gamma_{d+1}F_2$.  The program first reduces the words
without truncation and then independently verifies their matrices
of dimensions $3,7,15,31$ in the successive Magnus modules.
For the final assertion, in the class-two quotient of $B_3$ put
$k=[a,b]$, which is central.  The braid relation implies that the
abelianized images of $a,b$ agree; hence $b=a u$ with $u$ in the
central commutator subgroup, so $[a,b]=1$.  The relation then gives
$a=b$.  Thus $\gamma_2B_3/\gamma_3B_3=0$ and
$\gamma_2B_3=\gamma_3B_3$.  Inductively all later terms equal
$\gamma_2B_3$, proving that every nilpotent quotient factors through
$B_3^{\mathrm{ab}}=\mathbb Z$.  The full presentation proof is
therefore the argument in E.2, rather than the finite nilpotent test.
\end{proof}

\begin{statement}{D.8}{Proved}{The positive-power question over a finite local base}
Let $A=\mathbb Z/p^e$ and let $G/A$ be any group scheme, with no
linearity or affineness hypothesis.  If $g\in G(A)$ reduces to a
finite-order element of $G(\mathbb F_p)$, then $g$ has finite
order.  In particular, the positive-power conclusion holds for
all fibrewise quasi-unipotent sections over $\mathbb Z/4$.
\end{statement}
\begin{proof}
For a square-zero ideal $I\subset A$, the kernel of
$G(A)\to G(A/I)$ at the identity injects into the additive group
$\operatorname{Hom}(e^*\Omega^1_{G/A},I)$: the difference of two
ring maps with the same reduction is an $A$-derivation, and the
group law linearizes to addition.  This argument may be made in
an affine neighbourhood of the identity, since a section reducing
to the identity factors through that neighbourhood.  Apply it
successively to $A/p^j\to A/p^{j-1}$.  Each kernel has exponent
dividing $p$, so the kernel of $G(A)\to G(\mathbb F_p)$ is killed
by some $p$-power.  A power of $g$ lies in that kernel, and a
further power is the identity.  The constant homomorphism from
$\mathbb G_a$ has that value.  In characteristic $p$, the
manuscript's definition of quasi-unipotence implies finite order,
because a value on $\mathbb G_a$ is killed by $p$.  This proves
the last assertion and removes this particular non-linear
test case from the unresolved range of LM1 E.11.
\end{proof}

\begin{statement}{D.9}{Proved}{Scope of the integral-model questions}
The positive-power theorem holds for separated finitely presented
group schemes over every quasi-compact quasi-separated base
(D.17--D.21). Quasi-compactness alone is insufficient (D.14).
The smooth-model existence question at \S52.3, folio 754,
scan 213, concerns the preceding adelic factorization and
character conditions (45)--(49); it has the affirmative form
proved in D.22. The stronger integral-point recap on scan 213
and uniqueness from those data alone are corrected in D.23.
The later conditions (55)--(64) are supplementary hypotheses
and a definition, not an assertion that every Lie approximation
admits a model satisfying them all.
\end{statement}
\begin{proof}
The mathematical assertions are proved in D.14 and D.17--D.23
below. The distinction of quantifiers is directly visible in
the original images. Scan 213 requests an existence theorem
for the preceding conditions; scan 215 explicitly declares
primitive factorization (55) a supplementary hypothesis.
Scans 216--219 introduce further desired conditions, and
scan 219 defines an integral nondegenerate approximation when
they hold. D.4 and D.10 refute two proposed automatic consequences;
they do not make the remaining definitional hypotheses into
unproved universal existence claims. The precise discrepancy
between $\mathbb Z$ and $\widehat{\mathbb Z}$ in the recap
is recorded with both readings in D.23 and the reading ledger.
\end{proof}

\begin{statement}{D.10}{Refuted}{The discrete relations and integer gluing do not present the integral sheaf}
The presentation proposed as a question at scan 217, equation (57),
does not suffice in the category of group sheaves on the big etale
site.  Even for $H=\mathbb G_a^2$, $\Pi_0=\mathbb Z^2$, and the
two primitive coordinate one-parameter subgroups, the relations
of $\Pi_0$ and the two identifications on integer points do not
force the two additive subgroups to commute.
\end{statement}
\begin{proof}
Fix a prime $p$.  On $\mathbb F_p$-schemes let $U$ be the Heisenberg
group with coordinates $(a,b,c)$ and law
\[
 (a,b,c)(a',b',c')=(a+a',b+b',c+c'+ab').
\]
Regard $T\mapsto U(T\times_{\mathbb Z}\mathbb F_p)$ as a group
sheaf on the big etale site of $\operatorname{Spec}\mathbb Z$.
There are two homomorphisms from $\mathbb G_a$ to this sheaf,
\[
 t\longmapsto(t^p-t,0,0),\qquad
 s\longmapsto(0,s^p-s,0).
\]
The polynomial $t^p-t$ is additive in characteristic $p$.  Both
maps vanish on every integer section, so together with the zero
homomorphism from the constant group $\mathbb Z^2$ they satisfy
all the proposed relations.  Their commutator over
$\mathbb F_p[t,s]$ is $(0,0,(t^p-t)(s^p-s))\ne1$.
Consequently they cannot extend to a homomorphism from
$\mathbb G_a^2$, whose two coordinate subgroups commute.  This
contradicts the asserted universal property of the proposed
presentation.  The obstruction survives etale localization,
because a nonzero section of this represented sheaf cannot
become zero under a faithfully flat cover.
\end{proof}

The constant component groups and open adelic images in scans
217--219, equations (58)--(64), are expressly additional hypotheses,
not consequences of quasi-unipotence.  The associated right-exact
component sequences hold for smooth surjections with connected
identity fibres; they do not create those hypotheses.  The
finite-index subgroup discussed at scan 220 is not specified by
the definition.  Scan 221 begins a new construction from a chosen
representation and density assumption.  These choices, and its
possible continuation to the adelic group, belong to the category
of admissible targets in Part B rather than to D.1's uniqueness
of a single rational parameter.

\begin{statement}{D.11}{Proved}{The nonlinear dual-number case}
Let $k$ be a field of characteristic zero,
$A=k[\epsilon]/(\epsilon^2)$, and $G$ a separated
finite-presentation group scheme over $A$, not assumed
affine, flat or linear. If $g\in G(A)$ reduces to a
unipotent element of $G_k(k)$, there is a unique
$A$-group homomorphism $\mathbb G_{a,A}\to G$ taking $1$ to $g$.
If its reduction is only quasi-unipotent, the same
assertion holds for some positive power of $g$.
Here unipotent has its intrinsic field meaning:
the element belongs to a unipotent algebraic subgroup.
In particular this settles the nonlinear
$\Q[\epsilon]/(\epsilon^2)$ test left in LM1.
\end{statement}
\begin{proof}
We construct a substitute for a linear representation.
Let $R=\operatorname{Res}_{A/k}G$. It is a separated
finite-type $k$-group scheme. Indeed restriction of
scalars of an affine finite-presentation $A$-scheme
is affine of finite type by writing each coordinate
as $u+\epsilon v$ and equating the two coefficients
in each defining equation. For an arbitrary affine
open $W\subset G$, its restriction is an open
subfunctor of $R$: its complement is the image of
the closed complement under the finite projection
$\Spec A\times T\to T$. These open subfunctors cover,
because that projection is a universal homeomorphism
and hence each geometric fibre has only one support
point, which lies in an affine open of $G$.
They glue to the asserted finite-type scheme.
This is also the nilpotent-base special case of
\citet[Proposition 97.11.5, Tag 05Y8]{D_StacksRestriction}.
Restriction of scalars preserves products, so the
group law glues as well.

There is a reduction homomorphism $r:R\to G_k$.
Its kernel is the vector group
\[
 K=\Lie(G_k)\otimes_k(\epsilon k).
\]
To verify this even when $G$ is not flat, a point
of the kernel over a $k$-algebra $B$ is a lift of
the identity across the square-zero ideal
$\epsilon B\subset B[\epsilon]/\epsilon^2$.
Subtracting the given identity section identifies
such lifts with the derivations from its relative
cotangent space to $\epsilon B$. Base change of
relative differentials identifies that cotangent
space with the cotangent space of $G_k$; the group
law on these derivations is addition. This proves
the formula functorially.

We do not assume that $r$ is surjective. Every
finite-type group scheme over a characteristic-zero
field is smooth (Cartier's theorem,
\citealp[Lemma 39.8.2, Tag 047N]{D_StacksCartier}),
so the usual algebraic-group image theorem applies.
Let $H$ be the reduced Zariski closure of the
cyclic subgroup generated by $r(g)$. Since $r(g)$
is unipotent, $H$ is either the trivial group or
a copy of $\Ga$, by the finite logarithm in a
unipotent subgroup. The closed image of $r$
contains this cyclic subgroup and hence contains
$H$. Thus
\[
 1\longrightarrow K\longrightarrow r^{-1}(H)
           \longrightarrow H\longrightarrow1
\]
is an fppf exact sequence. The middle group is
affine: its map to the affine scheme $H$ is a
vector-group torsor, and such torsors over an
affine scheme are trivial, since their class is
in $H^1(H,\mathcal O_H^{\dim K})=0$.
It is therefore a unipotent group, since both
kernel and quotient are unipotent
(\citealp[Corollary 14.7]{D_Milne2022}).

The exponential isomorphism for a characteristic-zero
unipotent group
(\citealp[Proposition 14.32]{D_Milne2022}) now gives
$t\mapsto\exp(t\log g)$ in $r^{-1}(H)\subset R$.
By the defining adjunction of restriction of scalars
this is an $A$-group homomorphism $\mathbb G_{a,A}\to G$
whose value at $1$ is $g$.
Two such maps correspond to $k$-morphisms
$\mathbb G_{a,k}\to R$ agreeing at every integer, since both
take $n$ to $g^n$. The integers are Zariski dense
in $\mathbb G_{a,k}$, and $R$ is separated; the maps agree.
For quasi-unipotent reduction first take a power
whose reduction is unipotent and apply the result.

The nonflat test $G=\Spec A[t]/(\epsilon t)$
illustrates why the image qualification matters:
$R\simeq\mathbb G_{a,k}$, its reduction map has zero image,
and every $g=\epsilon a$ is nevertheless reached
by $t\mapsto\epsilon at$. Thus flatness has not
been silently inserted into the argument.
\end{proof}
